\begin{remarks}[2.7.3]
\label{IV.2.7.3-fr}
\begin{enumerate}
\item[(i)] De (2.6.1), (2.6.4) y (2.5.4.1) resulta que las conclusiones
de (2.7.1) siguen siendo válidas en los casos (i), (iii), (v), (vii) y
(xvi) cuando solo se supone que $g$ es cuasicompacto y
\emph{cuasifielmente plano} (2.3.3); ya se ha observado que (2.7.1) es válida
en el caso (ii) cuando solo se supone que $g$ es sobreyectivo y cuasicompacto.

\item[(ii)] Con las notaciones y las hipótesis de (2.7.1), puede ocurrir
que \emph{$f$ sea propio y $f'$ proyectivo sin que $f$ sea cuasiproyectivo}. En
efecto, Hironaka [34] dio un ejemplo de morfismo propio no proyectivo
$f:X\to Y$, donde $X$ e $Y$ son dos esquemas algebraicos regulares
$(\mathbf{0}_{\mathrm I},4.1.4)$ sobre un mismo cuerpo $k$, siendo $Y$ proyectivo sobre $k$;
además, $Y$ es reunión de dos abiertos afines $Y_i$ $(i=1,2)$ tales que
$f_i:X\times_Y Y_i\to Y_i$ es proyectivo para $i=1,2$. Sea entonces
$Y'=Y_1\amalg Y_2$ el preesquema suma; es claro que el morfismo
canónico $g:Y'\to Y$, cuyas restricciones a $Y_1$ y $Y_2$ son las
inyecciones canónicas, es fielmente plano, y es cuasicompacto en virtud de
(\textbf I, 5.5.10); sin embargo, aunque $f':X\times_Y Y'\to Y'$ (que coincide con
$f_i$ en cada uno de los $Y_i$) es proyectivo (\textbf{II}, 5.5.6), no ocurre lo
mismo con $f$. Existe, por tanto, un $\mathcal O_{X'}$-módulo invertible
$\mathcal L'$ que es amplio para $f'$ pero que \emph{no es de la forma
$g'^*(\mathcal L)$ para un $\mathcal O_X$-módulo invertible $\mathcal L$}, en
virtud de (2.7.2).

\item[(iii)] Bajo las hipótesis de (2.7.1), puede ocurrir que $f'$ sea
un \emph{isomorfismo local} sin que $f$ sea una inmersión local. Sean, en
efecto, $k$ un cuerpo, $\overline k$ una clausura algebraica de $k$, $K$ una
extensión separable de grado finito de $k$, distinta de $k$; entonces el
morfismo estructural $f:X\to Y$, donde $X=\operatorname{Spec}(K)$ e
$Y=\operatorname{Spec}(k)$, no es una inmersión local, pero, si se toma
$Y'=\operatorname{Spec}(\overline k)$, el morfismo $Y'\to Y$ es
fielmente plano y cuasicompacto y $f'$ es un isomorfismo local, puesto que
$X'=X\times_Y Y'$ es suma de un número finito de esquemas isomorfos a $Y'$.
\end{enumerate}
\end{remarks}

